\section*{第 \ref{ch:g6:measure} 章——周长、面积、体积}

\begin{solution}{exo:g6:measure:1}
$3.5$ m $= 350$ cm; \quad $420$ mm $= 42$ cm; \quad
$0.8$ km $= 800$ m; \quad $25\,000$ m $= 25$ km.
\end{solution}

\begin{solution}{exo:g6:measure:2}
长方形：$2 \times (8 + 3.5) = 2 \times 11.5 = 23$ cm。
正方形：$4 \times 6.2 = 24.8$ cm。
三角形：$5 + 7 + 9 = 21$ cm。
\end{solution}

\begin{solution}{exo:g6:measure:3}
一圈：$2\pi \times 50 = 100\pi \approx 314$ m。对于 $2$ km
$= 2000$ m 时：$2000 \div 314 \approx 6.4$，故需要 $7$ 整圈
（$6$ 圈只能走约 $1\,885$ m）。
\end{solution}

\begin{solution}{exo:g6:measure:4}
长方形：$9 \times 4 = 36$ cm$^2$。正方形：$7^2 = 49$ m$^2$。
\omterm{def:g6:shapes:triangles}{直角三角形}：$\frac{6 \times 10}{2} = 30$ cm$^2$。
圆：$\pi \times 3^2 = 9\pi \approx 28$ cm$^2$。
\end{solution}

\begin{solution}{exo:g6:measure:5}
$3$ m$^2 = 30\,000$ cm$^2$; \quad
$45\,000$ cm$^2 = 4.5$ m$^2$; \quad
$2.5$ L $= 2\,500$ cm$^3$; \quad
$4\,500$ L $= 4.5$ m$^3$.
\end{solution}

\begin{solution}{exo:g6:measure:6}
围栏（\omterm{def:g6:measure:perimeter}{周长}）：$2 \times (120 + 85) = 2 \times 205 = 410$ m。
\omterm{def:g6:measure:area}{面积}：$120 \times 85 = 10\,200$ m$^2$。
\end{solution}

\begin{solution}{exo:g6:measure:7}
\omterm{def:g5:solids:def}{长方体}：$5 \times 4 \times 10 = 200$ cm$^3$。\omterm{def:g5:solids:def}{正方体}：
$3 \times 3 \times 3 = 27$ cm$^3$。
\end{solution}

\begin{solution}{exo:g6:measure:8}
\omterm{def:g6:measure:perimeter}{周长} $16$ cm 意味着长 $+$ 宽 $= 8$ cm。例如
$6 \times 2$（\omterm{def:g6:measure:area}{面积} $12$ cm$^2$）和 $5 \times 3$（\omterm{def:g6:measure:area}{面积} $15$ cm$^2$）
——或者正方形 $4 \times 4$（\omterm{def:g6:measure:area}{面积} $16$ cm$^2$）。越接近
正方形，\omterm{def:g6:measure:area}{面积}就越大。
\end{solution}

\begin{solution}{exo:g6:measure:9}
\omterm{def:g6:measure:area}{面积}：$10 \times 2 + 2 \times 6 = 20 + 12 = 32$ cm$^2$。

\omterm{def:g6:measure:perimeter}{周长}（竖杆位于横杆正下方中央）：沿着图形走一圈，
\[
10 + 2 + 4 + 6 + 2 + 6 + 4 + 2 = 36 \text{ cm}
\]
（上边；横杆右端；横杆下侧右边；竖杆右边；竖杆
底端；竖杆左边；横杆下侧左边；横杆左端）。
\end{solution}

\begin{solution}{exo:g6:measure:10}
\emph{1.} $V = 25 \times 10 \times 2 = 500$ m$^3$.

\emph{2.} $500$ m$^3 = 500 \times 1000 = 500\,000$ L.

\emph{3.} $500\,000 \div 50\,000 = 10$ 小时。
\end{solution}

\begin{solution}{exo:g6:measure:11}
花园：$20^2 = 400$ m$^2$。池塘：$\pi \times 3^2 = 9\pi$ m$^2$。
草地：$400 - 9\pi \approx 400 - 28.3 \approx 372$ m$^2$。
\end{solution}

\begin{solution}{exo:g6:measure:12}
\emph{1.} 小条：$15 \times 6 \times 1 = 90$ cm$^3$。大条：
$30 \times 12 \times 2 = 720$ cm$^3$。\omterm{def:g6:measure:volume}{体积}变为原来的
$720 \div 90 = 8$ 倍——三个尺寸各加倍，\omterm{def:g6:measure:volume}{体积}就变为原来的
$2 \times 2 \times 2 = 8$ 倍。

\emph{2.} 巧克力的量是八倍，价钱只是四倍：是的，
大条每克的性价比是小条的两倍。
\end{solution}

\begin{solution}{pb:g6:measure:1}
\textbf{1.} $1 \times 9$：\omterm{def:g6:measure:perimeter}{周长} $2 \times (1 + 9) = 20$ m，
\omterm{def:g6:measure:area}{面积} $9$ m$^2$。\ $2 \times 8$：\omterm{def:g6:measure:perimeter}{周长} $2 \times 10 = 20$ m，
\omterm{def:g6:measure:area}{面积} $16$ m$^2$。

\textbf{2.} \omterm{def:g6:measure:perimeter}{周长}是（长 $+$ 宽）的两倍
（见\cref{ex:g6:measure:perimeter}），因此长 $+$ 宽
$= 20 \div 2 = 10$ m。列表如下：
\begin{center}
\renewcommand{\arraystretch}{1.2}
\begin{tabular}{c|ccccc}
围栏 & $1 \times 9$ & $2 \times 8$ & $3 \times 7$ & $4 \times 6$
& $5 \times 5$ \\
\hline
\omterm{def:g6:measure:area}{面积}（m$^2$） & $9$ & $16$ & $21$ & $24$ & $25$ \\
\end{tabular}
\end{center}
正方形 $5 \times 5$ 最好。

\textbf{3.} $4.5 \times 5.5 = 24.75$ 与 $4.9 \times 5.1 =
24.99$：越来越接近 $25$，但始终小于它。猜想：
正方形胜过\omterm{def:g6:measure:perimeter}{周长}为 $20$ 的\emph{每一个}长方形，
包括边长为\omterm{def:g6:decimals:places}{小数}的情形。

\textbf{4.} 由表可知：和为
$10$ 的两个数的积，在两数\emph{相等}时最大（$5$ 和
$5$）。一般地，和固定时，积在两部分
相等时最大——两部分越不相等，
积就越小。

\textbf{5.} 从 $5 \times 5$ 正方形中去掉最上面一行
——即由 $5$ 个单位正方形组成的一条窄条——余下 $5 \times 4$ 的围栏；
再在一端接上一列 $4$ 个单位正方形，使
围栏变为 $6 \times 4$。共拿走五个单位正方形，
只补回四个：所接的窄条沿 $4$ 边放置，
比被去掉的窄条所覆盖的 $5$ 边短。净损失：
一平方米（$25 \to 24$）。下一步，从
$6 \times 4$ 变为 $7 \times 3$，以一行 $6$ 个
换成一列 $3$ 个，又损失三平方米（$24 \to 21$）：围栏离正方形越远，
放弃的窄条就越长，
补回的窄条就越短。

\textbf{6.} 由 $w + L + w = 20$，即 $L = 20 - 2w$：
\begin{center}
\renewcommand{\arraystretch}{1.2}
\begin{tabular}{c|ccccccc}
$w$ & $1$ & $2$ & $3$ & $4$ & $5$ & $6$ & $7$ \\
\hline
$L$ & $18$ & $16$ & $14$ & $12$ & $10$ & $8$ & $6$ \\
\omterm{def:g6:measure:area}{面积} & $18$ & $32$ & $42$ & $48$ & $50$ & $48$ & $42$ \\
\end{tabular}
\end{center}
胜出者是 $5 \times 10$，\omterm{def:g6:measure:area}{面积}为 $50$ m$^2$——沿墙方向的长
恰为宽的两倍：\emph{不是}正方形。

\textbf{7.} 将围栏关于墙作对称：围栏加上它的镜像
构成一个加倍的围栏，篱笆用了两次：共 $40$ m
篱笆围住四周（墙的一侧此时位于内部）。由第一部分可知，
\omterm{def:g6:measure:perimeter}{周长}为 $40$ 的最优长方形是正方形，$10 \times
10$，\omterm{def:g6:measure:area}{面积}为 $100$ m$^2$。实际最优围栏是这个最优加倍围栏的\omterm{def:g3:division:half}{一半}：
沿墙 $10$ m，进深 $5$ m，\omterm{def:g6:measure:area}{面积}
$50$ m$^2$——恰好就是表中的胜者，现在得到了解释。

\textbf{8.} \omterm{def:g6:measure:perimeter}{周长} $20 = 2 \times \pi \times r$，故
$r = 20 \div 6.28 \approx 3.18$ m。\omterm{def:g6:measure:area}{面积}：
$\pi r^2 \approx 3.14 \times 3.18 \times 3.18 \approx
31.8$ m$^2$——明显大于正方形的 $25$ m$^2$。
狄多深知自己在做什么：对于给定长度的边界，
圆所围的\omterm{def:g6:measure:area}{面积}最大。

\textbf{9.} \omterm{def:g6:measure:perimeter}{周长}：$1 \times 36$：$74$ m；\ $2 \times 18$：
$40$ m；\ $3 \times 12$：$30$ m；\ $4 \times 9$：$26$ m；\
$6 \times 6$：$24$ m。正方形再次胜出，给定\omterm{def:g6:measure:area}{面积}下
所需篱笆最少。这恰是第一部分的反过来：\omterm{def:g6:measure:perimeter}{周长}固定时\omterm{def:g6:measure:area}{面积}最大，
或\omterm{def:g6:measure:area}{面积}固定时\omterm{def:g6:measure:perimeter}{周长}最小：两者都让正方形夺冠
（在一切长方形之外，还有圆）。

\textbf{10.} 圆形墙是所围空间的最短边界（见第 8 题），因此
圆形的罐、管或水箱
用料最少而容量最大——材料就是金钱。

\textbf{11.} 例如 $50$ m $\times$ $0.01$ m：\omterm{def:g6:measure:perimeter}{周长}
$2 \times (50 + 0.01) = 100.02$ m，\omterm{def:g6:measure:area}{面积}
$50 \times 0.01 = 0.5$ m$^2$。又长又细：数公里的
边界，几乎没有土地。

\textbf{12.} 错误。第 11 题的围栏\omterm{def:g6:measure:perimeter}{周长}
为 $100.02$ m，\omterm{def:g6:measure:area}{面积}为 $0.5$ m$^2$，而 $5 \times 5$ 的围栏
\omterm{def:g6:measure:perimeter}{周长}为 $20$ m，\omterm{def:g6:measure:area}{面积}为 $25$ m$^2$：\omterm{def:g6:measure:perimeter}{周长}更大，
而\omterm{def:g6:measure:area}{面积}小得多。

\textbf{13.} \omterm{def:g6:measure:area}{面积}：缺少一个单位正方形，
$24 - 1 = 23$ cm$^2$。\omterm{def:g6:measure:perimeter}{周长}：沿边走一圈，缺口处
以小正方形的三条边代替一段 $1$ cm 长的直墙，
即 $1 + 1 + 1 = 3$ cm，\omterm{def:g6:measure:perimeter}{周长}由 $20$ cm 增至
$20 - 1 + 3 = 22$ cm。去掉材料使
边界\emph{变长}。

\textbf{14.} 对真实海岸的每一次放大，都会显露出新的海湾、礁石
和小河汊——缺口叠着缺口，第 13 题表明
每个缺口都增加边界长度，而\omterm{def:g6:measure:area}{面积}几乎不变。
因此，所测长度随细节层次的提高不断增大，
（“海岸线悖论”），而所测\omterm{def:g6:measure:area}{面积}趋于稳定：
\omterm{def:g6:measure:perimeter}{周长}与\omterm{def:g6:measure:area}{面积}确实是彼此独立的量。

\textbf{15.} 对称论证：加倍后的围栏需用
$2 \times 36 = 72$ m 篱笆，\omterm{def:g6:measure:perimeter}{周长}为
$72$ 的最优长方形是 $18 \times 18$ 的正方形。因此最优围栏为
沿墙 $18$ m，进深 $9$ m。篱笆检验：
$9 + 18 + 9 = 36$ m，\omterm{def:g6:measure:area}{面积} $18 \times 9 = 162$ m$^2$。在开阔的
田野上，最优的是 $9 \times 9$ 的正方形，\omterm{def:g6:measure:area}{面积}为 $81$ m$^2$。
谷仓的墙恰好使农夫的土地\emph{翻倍}。
\end{solution}
